\chapter{Choisir}\label{ch:in:choosing}

Quatre offres, un poste de quant en banque, un poste de trading chez un teneur de marché, un poste d'analyste sur une plateforme et un poste d'ingénieur
dans une entreprise technologique, diffèrent par leur rémunération espérée, par leur fourchette, par ce que prélèvent l'impôt et le logement, par les horaires, par ce que les cinq
années suivantes enseignent et par la probabilité que le poste dure. Aucun chiffre unique ne les classe~: ce sont les pondérations de chaque personne qui le font, et ces pondérations sont
difficiles à énoncer. Les questions utiles sont de savoir quelles offres aucune pondération ne peut sauver, à quelle fréquence chaque offre arrive en tête quand les pondérations sont incertaines,
et de combien une pondération doit bouger pour que le choix change. Ce chapitre construit cette analyse
à partir des outils des chapitres précédents et clôt le livre.

\section{Les attributs qui différencient les types d'entreprise}

Les chapitres des parties~II et~III ont mesuré ce qui diffère.
\begin{itemize}
\item \textbf{Rémunération}, sa structure et sa fourchette (chapitres~13 et~14)~: \omterm{def:in:how-pay-works:base}{salaire de base} et prime, report et \omterm{def:in:how-pay-works:rsu}{déchéance},
rémunérations à la formule, et les fourchettes qu'affichent les dépôts réglementaires.
\item \textbf{Lieu} (chapitres~15 et~27)~: impôt, logement et niveaux de prix.
\item \textbf{Horaires} (chapitre~26)~: la couverture qu'exige un desk et les astreintes qu'il partage.
\item \textbf{Le travail} (chapitres~16 à~25)~: ce qu'un poste produit, à quelle vitesse on apprend si l'on réussit (la \omterm{def:in:quant-researcher:feedback}{vitesse du retour d'information} du chapitre~17) et
où il mène (chapitre~28).
\item \textbf{Sécurité}~: les mécanismes d'arrêt d'une plateforme (chapitre~22) et la probabilité
qu'un emploi prenne fin pour des raisons que son titulaire ne contrôle pas.
\end{itemize}

\begin{definition}[Modèle de valeur multi-attributs, pondération par écarts (swing weighting)]\label{def:in:choosing:mavm}
Un \emph{modèle de valeur multi-attributs}\index{modèle de valeur multi-attributs} (multi-attribute value model) note chaque option sur chaque attribut,
ramène chaque attribut de la pire à la meilleure option, puis additionne les notes ainsi ramenées à l'échelle, affectées de pondérations (Keeney et Raiffa).
La \emph{pondération par écarts}\index{pondération par écarts} (swing weighting) fixe les pondérations en demandant combien le décideur valorise le passage de chaque attribut de son pire à son meilleur niveau,
relativement aux autres, plutôt que de demander à quel point l'attribut compte
dans l'absolu.
\end{definition}

\section{Le risque de rémunération comme équivalent certain}

Le chapitre~13 a transformé la distribution d'une rémunération globale en équivalent certain~: le montant sûr qui procure la même utilité espérée
sous une aversion relative au risque constante (livre~4, chapitre~9). Avec une aversion relative au risque de 3 et 500\,000~\$ d'autres actifs,
les quatre offres du chapitre valent 1\,186\,354~\$ (quant en banque), 1\,365\,186~\$ (teneur de marché), 903\,603~\$ (analyste de plateforme) et 845\,964~\$ (ingénieur technologique) sur cinq ans.
L'équivalent certain intègre le risque dans la rémunération~; le
modèle de valeur ci-dessous les sépare comme rémunération espérée et risque de rémunération (l'étendue entre le 10\up{e} et le 90\up{e}~centile
rapportée à la moyenne), de sorte que le poids du risque puisse être repéré et varié.

\begin{center}\footnotesize\setlength{\tabcolsep}{4pt}
\begin{tabular}{@{}>{\raggedright\arraybackslash}p{2.6cm}>{\centering\arraybackslash}p{1.8cm}>{\centering\arraybackslash}p{1.8cm}>{\centering\arraybackslash}p{1.8cm}>{\centering\arraybackslash}p{1.8cm}>{\centering\arraybackslash}p{1.8cm}>{\centering\arraybackslash}p{1.8cm}@{}}
\toprule
offre (ville) & rémuné\-ration espérée & risque de rémunération & impôt, logement & astreinte & appren\-tissage & sécurité\\
 & k\$, 5 ans & étendue / moyenne & part & nuits/mois & 1--5 & part\\
\midrule
quant en banque (Londres) & 1\,618 & 1.01 & 53.1~\% & 0.0 & 3 & 1.00\\
teneur de marché (Chicago) & 1\,949 & 1.16 & 40.8~\% & 3.8 & 4 & 1.00\\
analyste de plateforme (New York) & 1\,703 & 1.67 & 50.0~\% & 0.0 & 4 & 0.65\\
ingénieur technologique (Londres) & 1\,191 & 1.22 & 55.3~\% & 5.1 & 3 & 1.00\\
\bottomrule
\end{tabular}
\end{center}

Les chiffres de rémunération du tableau proviennent des offres du chapitre~13 et d'une quatrième construite de la même façon (le \omterm{def:in:how-pay-works:base}{salaire de base}
d'une entreprise technologique, auquel s'ajoute une attribution d'actions restreintes acquise sur quatre ans)~; chaque ensemble de rémunération est illustratif.
La part prélevée par l'impôt et le logement applique les outils du chapitre~27 à la rémunération annuelle moyenne de chaque offre dans sa ville~; les nuits d'astreinte utilisent les plannings du chapitre~26
(huit personnes pour le desk de teneur de marché en continu, six pour le service de l'ingénieur)~; l'apprentissage relève du jugement du décideur~;
la sécurité est la part des carrières de cinq ans dans lesquelles le modèle de rémunération ne couvre pas la perte d'emploi, ce que seul le registre
de la plateforme peut établir.

\section{Tempérament, compétences et le travail lui-même}

Certains attributs résistent aux chiffres et comptent le plus. La \omterm{def:in:quant-researcher:feedback}{vitesse du retour d'information} du chapitre~17 décide du délai au bout duquel une personne apprend si
elle est douée pour le poste~; la charge d'encadrement du chapitre~16 et les horaires du chapitre~26 décident de ce que ressemble une journée~; les mécanismes d'arrêt du chapitre~22
décident du coût d'une mauvaise année. Une personne qui n'aime pas la rémunération incertaine doit le dire dans la pondération du risque de rémunération~; une
personne qui veut construire des systèmes plutôt que de les trader doit le dire dans la note d'apprentissage. Le modèle ne connaît pas le tempérament~;
il oblige le décideur à l'énoncer et montre ce qui en découle.

\section{Un cadre multi-attributs et sa sensibilité}

\begin{definition}[Option dominée, acceptabilité de rang]\label{def:in:choosing:smaa}
Une alternative est une \emph{option dominée}\index{option dominée} lorsqu'une autre
est au moins aussi bonne sur chaque attribut et meilleure sur un~: aucune pondération ne peut la rendre la meilleure. L'\emph{acceptabilité de rang}\index{acceptabilité de rang} d'une alternative
pour un rang est la part des vecteurs de pondérations, tirés d'une distribution énoncée, sous lesquels l'alternative occupe ce rang
(analyse stochastique d'acceptabilité multicritère de Lahdelma, Hokkanen et Salminen).
\end{definition}

\begin{method}[Filtrage, pondération et analyse de sensibilité]\label{met:in:choosing:method}
Ramener chaque attribut à $[0,1]$ de la pire à la meilleure offre (en inversant là où moins vaut mieux). Écarter les offres dominées.
Avec les \omterm{def:in:choosing:mavm}{pondérations par écarts} $w$, la valeur de l'offre $i$ est $\sum_j w_j s_{ij}$. Tirer les pondérations uniformément au hasard sur le simplexe (une
loi de Dirichlet à paramètres unitaires) et compter combien de fois chaque offre occupe chaque rang. Pour
l'attribut d'intérêt, chercher le plus petit changement de son poids, les autres étant remis à l'échelle proportionnellement, qui modifie la meilleure offre.
\end{method}

Avec les points d'écart illustratifs (rémunération espérée 30, risque de rémunération 15, lieu 10, astreinte 15, apprentissage 20, sécurité 10), la
note de l'offre du teneur de marché est de 0.852, celle de la plateforme de 0.589, celle de la banque de 0.584 et celle de l'ingénieur technologique de 0.202.
L'offre technologique est dominée~: la banque paie davantage, avec une fourchette plus étroite, prélève moins au titre de l'impôt et du logement, n'exige
aucune astreinte et lui est égale sur l'apprentissage et la sécurité. Sous des pondérations tirées uniformément au hasard (la \cref{fig:in:choosing:ranks}), le
teneur de marché arrive en tête dans 79.5~\% des tirages, la banque dans 15.9~\%, la plateforme dans 4.6~\% et l'offre technologique jamais.

\begin{omfigure}
\begin{tikzpicture}
\begin{axis}[xbar stacked,width=10.5cm,height=5cm,xmin=0,xmax=100,ymin=0.4,ymax=4.6,bar width=12pt,ytick=data,
  yticklabels={{quant en banque},{teneur de marché},{analyste de plateforme},{ingénieur technologique}},yticklabel style={font=\footnotesize},
  y dir=reverse,xticklabel style={font=\footnotesize},xlabel={\small part des tirages de pondérations,~\%},table/col sep=comma,
  legend style={legend cell align=left,font=\scriptsize,at={(0.5,-0.3)},anchor=north,legend columns=4,/tikz/every even column/.append style={column sep=8pt}}]
\addplot[fill=omDef,draw=none] table[x=r1,y=pos]{figdata/industry/30-choosing/ranks.csv};
\addplot[fill=omDef!55,draw=none] table[x=r2,y=pos]{figdata/industry/30-choosing/ranks.csv};
\addplot[fill=omThm!45,draw=none] table[x=r3,y=pos]{figdata/industry/30-choosing/ranks.csv};
\addplot[fill=black!25,draw=none] table[x=r4,y=pos]{figdata/industry/30-choosing/ranks.csv};
\legend{premier,deuxième,troisième,quatrième}
\end{axis}
\end{tikzpicture}

\omcaption{\omterm{def:in:choosing:smaa}{Acceptabilité de rang} des quatre offres illustratives~: la part des 20\,000 vecteurs de pondérations, tirés uniformément sur
le simplexe des six attributs, sous lesquels chaque offre occupe chaque rang. Données~: \texttt{in\_choose.results}, issues de
\texttt{firm.careerdec.rank\_acceptability}.}\label{fig:in:choosing:ranks}
\end{omfigure}

L'avance du teneur de marché résiste au poids accordé au risque sur la rémunération~: ce poids doit passer
de 0.15 à 0.602, soit une hausse de 0.452, les autres poids étant réduits dans la même proportion, avant que l'offre de la banque ne devienne la meilleure (voir la \cref{fig:in:choosing:sweep}). Celui qui choisit
pour qui la dispersion de la rémunération pèse plus que la rémunération espérée, le lieu, les horaires, l'apprentissage et la sécurité réunis devrait retenir l'offre de la banque~; tout autre
au vu de ces chiffres, retient celle du teneur de marché.

\begin{omfigure}
\begin{tikzpicture}
\begin{axis}[width=10.5cm,height=5.6cm,xmin=0,xmax=0.95,ymin=0,ymax=1,xlabel={\small poids du risque sur la rémunération},
  ylabel={\small valeur},tick label style={font=\footnotesize},grid=major,grid style={gray!20},table/col sep=comma,
  legend style={font=\scriptsize,at={(0.5,-0.3)},anchor=north,legend columns=2,legend cell align=left,
  /tikz/every even column/.append style={column sep=8pt}}]
\addplot[omDef,thick] table[x=w,y=bank_quant]{figdata/industry/30-choosing/sweep.csv};
\addplot[omThm,thick] table[x=w,y=market_maker]{figdata/industry/30-choosing/sweep.csv};
\addplot[black!60,thick,dashed] table[x=w,y=platform_analyst]{figdata/industry/30-choosing/sweep.csv};
\addplot[black!30,thick] table[x=w,y=technology_engineer]{figdata/industry/30-choosing/sweep.csv};
\draw[black!50,dashed] (axis cs:0.15,0) -- (axis cs:0.15,1);
\draw[black!50,dashed] (axis cs:0.602,0) -- (axis cs:0.602,1);
\legend{quant de banque,teneur de marché,analyste de plateforme,ingénieur technologique}
\end{axis}
\end{tikzpicture}

\omcaption{Valeur de chaque offre à mesure que le poids du risque sur la rémunération varie, les autres poids étant remis à l'échelle proportionnellement. Les traits pointillés
marquent le poids retenu (0.15) et le poids à partir duquel l'offre de la banque dépasse celle du teneur de marché (0.602). Données~:
\texttt{in\_choose.results} et \texttt{firm.careerdec.flip\_shift}.}\label{fig:in:choosing:sweep}
\end{omfigure}

\section{Revenir sur la décision}

Une décision sur un emploi se révise, et c'est dans la révision que surviennent la plupart des erreurs. Le livre~16, chapitre~26, décrit
les outils~: un journal de décisions tenu sur le moment, qui consigne les attributs, les poids et les anticipations, et protège ainsi contre le
biais du résultat (juger le choix d'après le déroulement de l'année plutôt que d'après ce qui était connu)~; et un pré-mortem, qui
demande, avant le changement, comment il pourrait échouer. Le modèle de mobilité du chapitre~28 indique quand partir coûte cher, et la \omterm{def:in:quant-researcher:feedback}{vitesse du retour d'information} du chapitre~17
indique combien de temps attendre avant de juger un poste de recherche. Le modèle de ce chapitre est le relevé de ce que le décideur
croyait~; sa valeur tient à ce qu'il soit recalculé lorsqu'une conviction change.

\section{Tutoriel~: quatre offres}\label{tut:in:choosing:tut}

\begin{tutorial}
\textbf{Objectif.} Noter quatre offres avec les outils des chapitres précédents, les filtrer, les pondérer et mesurer la robustesse du choix.
\textbf{État final~:} les \cref{fig:in:choosing:ranks,fig:in:choosing:sweep} et le tableau des attributs.
\begin{tutsteps}
\item \textbf{Attributs.} \texttt{firm.payoffer} pour la rémunération et sa dispersion~; \texttt{firm.locations} pour ce que prennent l'impôt et
le logement~; \texttt{firm.workload.oncall\_nights}~; des appréciations pour l'apprentissage~; les clôtures du modèle de rémunération pour la sécurité.
\item \textbf{Échelle et filtrage.} \texttt{firm.careerdec.scale} et \texttt{dominated}.
\item \textbf{\omterm{def:in:choosing:smaa}{Acceptabilité de rang}.} \texttt{rank\_acceptability} tire les poids sur le simplexe
(voir le \cref{lst:in:choosing:smaa}).
\omcode{code/firm/careerdec/firm_careerdec.py}{49}{55}{Acceptabilité de rang par tirage uniforme des poids sur le simplexe.}\label{lst:in:choosing:smaa}
\item \textbf{Sensibilité.} \texttt{flip\_shift} recherche la plus petite variation d'un poids qui modifie la meilleure offre
(voir le \cref{lst:in:choosing:flip}).
\omcode{code/firm/careerdec/firm_careerdec.py}{66}{78}{La plus petite variation d'un poids qui modifie le choix.}\label{lst:in:choosing:flip}
\end{tutsteps}
\end{tutorial}

Les résultats~: l'acceptabilité au premier rang de 79.5~\% (teneur de marché), 15.9~\% (quant de banque), 4.6~\% (analyste de plateforme), 0~\%
(ingénieur technologique)~; le poids du risque sur la rémunération doit augmenter de 0.452 pour modifier le choix.

\textbf{Pour aller plus loin.} Remplacer la distribution uniforme des poids par le \omterm{def:in:bank-markets-divisions:league}{classement} des attributs établi par le décideur (SMAA
ordinale)~; ajouter une corrélation entre les incertitudes des attributs~; relancer l'analyse dans un an, à la lumière de ce qui aura été appris.

\section{Construction~: firm.careerdec}\label{bld:in:choosing:careerdec}

\begin{build}
\bfield{Rôle} Choisir entre des offres à plusieurs attributs, et dire dans quelle mesure le choix résiste aux poids.
\bfield{Interface} \texttt{firm.careerdec}~: \texttt{scale(matrix, higher\_better)}~; \texttt{swing\_weights(points)}~;
\texttt{value(scaled, weights)}~; \texttt{dominated(matrix, higher\_better)}~;
\texttt{rank\_acceptability(scaled, n, rng)}~; \texttt{flip\_shift(scaled, weights, attr, step)}~; alimenté par \texttt{firm.payoffer}, \texttt{firm.aftertax},
\texttt{firm.locations}, \texttt{firm.workload} et \texttt{firm.careerpath}.
\bfield{Règles} Les échelles vont de la pire à la meilleure offre~; la dominance est vérifiée sur les valeurs brutes~; les poids totalisent un~; la
recherche de bascule modifie un poids et réajuste les autres proportionnellement.
\bfield{Tests d'acceptation} \texttt{code/firm/careerdec/tests/}~: mise à l'échelle et inversion~; une offre dominée est repérée~; deux
offres miroirs arrivent chacune premières la moitié du temps~; une bascule est trouvée là où elle doit l'être et jamais là où elle ne peut pas l'être.
\bfield{Pour aller plus loin} Poids ordinaux~; scores par intervalles~; une fonction valeur avec des interactions entre attributs.
\end{build}

\begin{omsources}
\item Keeney, R.~L. et H.~Raiffa (1976), \emph{Decisions with Multiple Objectives: Preferences and Value Tradeoffs},
Wiley (Cambridge University Press, 1993).
\item von Winterfeldt, D. et W.~Edwards (1986), \emph{Decision Analysis and Behavioral Research}, Cambridge University
Press.
\item Lahdelma, R., J.~Hokkanen et P.~Salminen (1998), SMAA---stochastic multiobjective acceptability analysis,
\emph{European Journal of Operational Research} 106(1), 137--143.
\item Chapitres~13, 15, 17, 22 et 26--28 de ce livre~; livre~4, chapitre~9~; livre~16, chapitre~26.
\end{omsources}

\section{Exercices}

\begin{exercise}[$\star$]\label{exo:in:choosing:1}
Déterminer quelle offre est dominée, et par laquelle.
\end{exercise}

\begin{exercise}[$\star$]\label{exo:in:choosing:2}
Calculer, en parts, les \omterm{def:in:choosing:mavm}{pondérations par écarts} (swing weighting) des six attributs.
\end{exercise}

\begin{exercise}[$\star$]\label{exo:in:choosing:3}
Déterminer quelle offre a le plus haut équivalent certain pour une aversion relative au risque égale à 3, et laquelle a la rémunération espérée la plus élevée.
\end{exercise}

\begin{exercise}[$\star\star$]\label{exo:in:choosing:4}
Expliquer pourquoi l'offre de la plateforme n'arrive premier que dans 4.6~\% des tirages de poids alors qu'elle est à égalité pour le meilleur score d'apprentissage.
\end{exercise}

\begin{exercise}[$\star\star$]\label{exo:in:choosing:5}
Expliquer quelle est la faiblesse d'un tirage uniforme des poids sur le simplexe.
\end{exercise}

\begin{exercise}[$\star\star$]\label{exo:in:choosing:6}
Un ami attribue à l'apprentissage 60~points de \omterm{def:in:choosing:mavm}{pondération par écarts}, les autres attributs étant pondérés comme dans le chapitre. Déterminer quelle offre le modèle retient.
\end{exercise}

\begin{exercise}[$\star\star\star$]\label{exo:in:choosing:7}
\emph{Programmation.} Retirer l'offre dominée et recalculer l'\omterm{def:in:choosing:smaa}{acceptabilité de rang} premier des trois autres.
\end{exercise}

\begin{exercise}[$\star\star\star$]\label{exo:in:choosing:8}
\emph{Trouver la faille.} «~L'offre du teneur de marché l'emporte dans 79.5~\% des pondérations, elle est donc le bon choix pour 79.5~\% des
personnes.~»
\end{exercise}

\section{Problème~: quatre offres}

\begin{problem}[{Problème du week-end --- quatre offres}]\label{pb:in:choosing:1}
Une jeune diplômée dispose des quatre offres du chapitre et demande laquelle accepter.

\medskip\noindent\textbf{Partie I --- Les attributs.}
\begin{enumerate}
\item Quels attributs diffèrent d'une offre à l'autre, et quels chapitres mesurent chacun d'eux~?
\item Définir le \omterm{def:in:choosing:mavm}{modèle de valeur multi-attributs} et la \omterm{def:in:choosing:mavm}{pondération par écarts} (swing weighting).
\item Énoncer les quatre équivalents certains pour une aversion relative au risque égale à 3.
\item Pourquoi garder séparées la rémunération espérée et le risque sur la rémunération~?
\item Quels attributs relèvent du jugement~?
\end{enumerate}

\medskip\noindent\textbf{Partie II --- Le modèle.}
\begin{enumerate}[resume]
\item Donner le tableau des attributs.
\item Définir une \omterm{def:in:choosing:smaa}{option dominée} et nommer l'offre dominée.
\item Énoncer les points de \omterm{def:in:choosing:mavm}{pondération par écarts} et les valeurs des offres.
\item Définir l'\omterm{def:in:choosing:smaa}{acceptabilité de rang} et énoncer la méthode.
\item Donner les \omterm{def:in:choosing:smaa}{acceptabilités de rang} premier.
\end{enumerate}

\medskip\noindent\textbf{Partie III --- Sensibilité.}
\begin{enumerate}[resume]
\item De combien le poids du risque de rémunération doit-il varier pour modifier le choix~?
\item Que dit cela de l'offre de la banque~?
\item Qu'est-ce qui modifierait les scores d'apprentissage~?
\item Que changerait un versement plus élevé de la plateforme~?
\item Qu'ignore le modèle~?
\end{enumerate}

\medskip\noindent\textbf{Partie IV --- Le verdict.}
\begin{enumerate}[resume]
\item Énoncer le \emph{résultat nommé}~: l'indice d'\omterm{def:in:choosing:smaa}{acceptabilité de rang} premier de chaque offre sous des pondérations uniformes, et la plus petite
variation du poids du risque de rémunération qui modifie l'offre préférée.
\item Comment la diplômée doit-elle consigner la décision~?
\item Quand doit-elle la réexaminer~?
\item Qu'est-ce qui pourrait lui faire regretter cette décision, et comment distinguerait-elle la malchance d'un mauvais choix~?
\item En deux phrases, lui répondre.
\end{enumerate}
\end{problem}

\section{Questions d'entretien}

\begin{interviewq}[$\star$ \iqroles{chercheur, trader}]\label{iq:in:choosing:1}
Pourquoi souhaitez-vous ce poste plutôt que les autres que vous envisagez~?
\end{interviewq}

\begin{interviewq}[$\star$ \iqroles{chercheur}]\label{iq:in:choosing:2}
Comment compareriez-vous deux offres dont la rémunération diffère par sa structure mais non par sa valeur espérée~?
\end{interviewq}

\begin{interviewq}[$\star\star$ \iqroles{chercheur, développeur}]\label{iq:in:choosing:3}
Écrivez une fonction qui repère les \omterm{def:in:choosing:smaa}{options dominées}. Quelle est sa complexité~?
\end{interviewq}

\begin{interviewq}[$\star\star$ \iqroles{chercheur}]\label{iq:in:choosing:4}
Comment tireriez-vous uniformément des poids sur un simplexe~?
\end{interviewq}

\begin{interviewq}[$\star\star$ \iqroles{trader}]\label{iq:in:choosing:5}
Vous avez accepté un poste qui s'est mal passé la première année. Comment décidez-vous si c'était une mauvaise décision~?
\end{interviewq}

\begin{interviewq}[$\star\star\star$ \iqroles{chercheur}]\label{iq:in:choosing:6}
Deux attributs sont fortement corrélés entre vos options. Qu'est-ce que cela change pour un modèle de valeur additif, et que
feriez-vous à ce sujet~?
\end{interviewq}
